\scanpage{93}
- \mspage{IV-1012} -
L'implication a)$\Rightarrow$c) est immédiate \deleted{\illegible{1}} (\uline{0},15.1.13). Si c) est vérifiée, il résulte de (19.5.5) que $\psi_M$ est bijectif ; en outre,
\deleted{\unsure{$\operatorname{gr}_{\underline K}^\bullet(M)$} \illegible{6}}
$(\operatorname{gr}_{\underline K}^\bullet(M)\otimes_A(A/\underline J)=
=(\operatorname{gr}_{\underline K}^\bullet(M)\otimes_{A/\underline K}(A/\underline L)\unsure{)|/}$ est un $(A/\underline L)$-module plat, donc il en est de même de
$(\operatorname{gr}_{\underline K}^\bullet(M)\otimes_A(A/\underline J))[T_1,\ldots,T_n]$, et par suite de $\operatorname{gr}_{\underline L}^\bullet(M)$ puisque $\psi_M$ est un $(A/\underline L)$-isomorphisme ; donc c) entraîne b). Le fait que a) implique b) montre aussitôt que a) implique aussi d).
\margin{/(Bourbaki, \uline{Alg.}, chap. II, $3^\mathrm{e}$ éd., § 3, n° 7, cor. 2 de la prop. 6)}

Lorsque tout quotient d'un $\operatorname{gr}_{\underline K}^p(M)$ est idéalement séparé pour $\underline J$, il résulte de (19.5.5) que la condition b) entraîne que $f$ est $(\operatorname{gr}_{\underline K}^\bullet(M))$-régulière ; en outre, comme
$(\operatorname{gr}_{\underline K}^\bullet(M)\otimes_{A/\underline K}(A/\underline L))[T_1,\ldots,T_n]$, isomorphe à $\operatorname{gr}_{\underline L}^\bullet(M)$, est alors un $(A/\underline L)$-module plat, \deleted{\unsure{et}} il en est \deleted{\illegible{1}} de même de
$\operatorname{gr}_{\underline K}^\bullet(M)\otimes_{A/\underline K}(A/\underline L)$, qui en est un facteur direct ; cela prouve donc que b) entraîne alors c). D'autre part, c) entraîne alors a) en vertu de (\uline{0},15.1.21) (où on peut remplacer l'hypothèse noethérienne par l'hypothèse que $M$ est idéalement séparé pour $\underline J$, en vertu de $(\underline{0_{\mathrm{III}}},10.2.1)$).

Il reste donc à prouver que lorsque $A$ est noethérien, $M$ de type fini et que les \deleted{\illegible{1}} $f_i$ appartiennent au radical de $A$, d) entraîne a). Il existe par hypothèse un $A$-module \uline{libre de type fini} \inserted{$N$} et une suite exacte
$0\longrightarrow R\longrightarrow N\longrightarrow\unsure{M}\longrightarrow0$.
Il suffit de prouver que d) entraîne b), puisque tout quotient d'un sous-module d'un $\operatorname{gr}_{\underline K}^p(M)$ est alors idéalement séparé pour $\underline J$, en vertu du fait que $\operatorname{gr}_{\underline K}^p(M)$ est un $A$-module de type fini \deleted{\illegible{1}}, que $\underline J$ est contenu dans le radical de $A$, \deleted{\illegible{1}} et que $A$ est noethérien (Bourbaki, \uuline{Alg. comm.}, chap. III, § 5, n° 1) ; en d'autres termes, il s'agit de voir que $\psi_M$ est \uline{bijectif}.
\deleted{\illegible{5}}
\mspage{92}
